% Consolidation only. No preamble. Requires amsmath, amssymb and hyperref.
% Saved audits, not new research; 7 October 2026.
\section{Boolean prime ideals, compactness, and model variations}
\label{sec:bpi-variants}

This section records consequences of the released OpenAI construction
and adaptations checked in the present research. The construction itself
is credited to OpenAI, at commit
\href{https://github.com/openai/math/commit/adc7f1241b42e322a6451854ab7e4b4c146bf78a}{\texttt{adc7f124}}.
The new arguments below have mathematical drafts and separate internal
AI-agent audits. That is distinct from external peer review and from a
formal proof of their full statements. Their novelty has not been
certified. Older working notes saying that BPI in this model is
``undecided'' are superseded by the six-cycle theorem below.

Write $X\simeq Y$ for equipotence, and $X\preccurlyeq Y$ for the existence
of an injection. The relevant established package in the original model
$W$ is
\[
 \mathsf{ZF}+\mathsf{PP}+\mathsf{AC}_{\rm WO}+\mathsf{DC}
 +\mathsf{FB}+\mathsf{NDS}+\neg\mathsf{AC}.
\]
Here $\mathsf{FB}$ is the finite cardinal basis property defined earlier,
not an abbreviation for finite choice. The additional structural inputs
are $\mathsf{SVC}(A)$, the quotient dichotomy $Q(A)$, and the least
non-well-orderable cardinal $|A|$. Every set embeds into $A\times\eta$
for an ordinal $\eta$, every non-well-orderable set contains a copy of
$A$, and $A^{<\omega}\simeq A^2\simeq A$. The right action of the ground
monoid $T$ on $A$ is free on each orbit. The native order is
\[
 a\leq_T b\quad\Longleftrightarrow\quad
 (\exists t\in T)\ a\cdot t=b.
\]
Left cancellation and conicality of $T$ make this a partial order.
Its graph is an internal hereditarily symmetric object; its definition
does not require treating the generic filter as an internal parameter.

\subsection{A partial order with no linear extension}

\textbf{Audited theorem.}
The native partial order on $A$ has no linear extension in $W$.
Consequently
\[
 W\models\neg\mathsf{OEP}+\neg\mathsf{BPI},
\]
where OEP says that every partial order has a linear extension and BPI
says that every nontrivial Boolean algebra has a prime ideal. The same
argument applies to the parameter variants described below. The
Ordering Principle still holds: every set can be linearly ordered,
but an arbitrary prescribed partial order need not be extended.

\emph{Proof sketch with the compatibility step.}
The source monoid has a distinguished $e\ne1$ and $\mu$ many probes
$m_\xi$ satisfying
\[
 m_\xi e=e,\qquad m_\xi\ne1,\qquad
 m_\xi T\cap m_\zeta T=eT\quad(\xi\ne\zeta).
\]
It also realizes every valid finite prescription of principal-cone
overlaps, with exactly the prescribed identifications. A diagram is a
ground injection $q:T\to I$; write $q_t(x)=q(tx)$ and $D_q=q[T]$.
Cohen conditions mention fewer than $\mu$ coordinate labels.

Suppose a name for a linear extension $L$ is supported by $D_p$ and
$(c_0,p)$ forces its asserted properties. Restrict $c_0$ to $D_p$.
Choose $q$ with $q_e=p$ and two probes $a,b$. The source local-decision
lemma lets $q$ decide their comparison after the Cohen generic is fixed.
A Cohen strengthening $d$, restricted to $D_q$, therefore forces, after
possibly exchanging $a,b$,
\[
 q\Vdash A_{q(a)}<_L A_{q(b)}.
\]
This fixes the diagram before choosing the following native step.

The cones $bm_\xi T\setminus eT$ are pairwise disjoint. Hence some
probe $t\ne1$ has
\[
 \operatorname{supp}_I(d)\cap q[btT\setminus eT]=\varnothing.
\]
Set $h=bt$. Then $he=e$, $aT\cap hT=eT$, and $b<_T h$.
On three copies, indexed modulo $3$, prescribe
\[
                   (i,hx)\sim(i+1,ax)\qquad(x\in T).
\]
This prescription is valid. Indeed, a two-step passage through the
middle copy must lie in $aT\cap hT=eT$; if $ax=hy=ez$, cancellation
and $ae=he=e$ give $x=y=ez$. Thus the indicated relation is already
transitive, its common root has identical parameters, and no copy
collapses onto itself.

Place these copies as diagrams $q_i$, all with $(q_i)_e=p$, and transport
$d$ to conditions $d_i$ by coordinate permutations fixing $D_p$.
The $d_i$ agree on the root. Off the root, each outgoing overlap is
silent in the corresponding $d_i$ by the choice of $t$. Their union
is therefore a valid Cohen condition. This is the essential step:
a bare swap of two diagrams would not guarantee compatible Cohen bits.
The diagram merger gives a common extension of the $q_i$. With
$x_i=A_{q_i(a)}$ and $y_i=A_{q_i(b)}$, it forces
\[
 x_i<_L y_i<_T x_{i+1}\quad(i\bmod3),
\]
and hence a strict six-cycle in $L$, a contradiction.

Forcing automorphisms here transport names and conditions. They are
not asserted to be automorphisms of the actual pure ZF universe.
The implication $\mathsf{BPI}\Rightarrow\mathsf{OEP}$ is classical
propositional compactness; see Jech, \emph{The Axiom of Choice} (1973),
Example 2.3.2, p.~19. Jech's \S7.2 already uses finite symmetry cycles
in models with OP and without OEP. The model-specific contribution here
is the silent probe cone and its exact Cohen-compatible merger.

\emph{Verification scope.}
The auxiliary \path{SixCycleAudit.lean} proves validity of this cyclic
overlap using the pinned source's actual
\texttt{OverlapPrescription.Valid}, exact root intersections, absence
of within-copy collapse, and compatibility of copied partial assignments.
Its compilation reported no \texttt{sorryAx}. It does \emph{not}
formalize the full symmetric-forcing argument or the conclusion
$W\models\neg\mathsf{BPI}$.
The principal records are
\path{NATIVE_ORDER_EXTENSION_SIX_CYCLE_PROOF.txt},
\path{INDEPENDENT_AUDIT_SIX_CYCLE_OEP_FAILURE.txt}, and
\path{SIX_CYCLE_LEAN_SCOPE.txt}.

\subsection{Explicit Boolean-algebra and graph witnesses}

\textbf{Boolean algebra.}
There is a nontrivial Boolean algebra $B_{\rm ord}$ of cardinality
exactly $|A|$ with no prime ideal. Use variables $p_{ab}$, $a,b\in A$,
and impose
\[
 \neg p_{aa},\quad p_{ab}\vee p_{ba}\ (a\ne b),\quad
 \neg(p_{ab}\wedge p_{ba})\ (a\ne b),\quad
 (p_{ab}\wedge p_{bc})\Rightarrow p_{ac},\quad
 p_{ab}\ (a<_T b).
\]
Take the free Boolean algebra modulo the ideal generated by the
negations of these constraints. Every finite subtheory is satisfiable:
linearly extend the native partial order on its finitely many mentioned
points. Thus the quotient is nontrivial. A prime ideal would give a
homomorphism to $2$, hence a forbidden native linear extension.

Finite terms have cardinality $A$, so $A$ surjects onto $B_{\rm ord}$.
Every well-orderable nontrivial Boolean algebra has a prime ideal in ZF:
recursively include the next enumerated element in a growing proper
ideal whenever this preserves properness. If it cannot be included,
its complement already belongs to the ideal. Limit unions stay proper.
Consequently $B_{\rm ord}$ is not well-orderable, and $Q(A)$ gives
$B_{\rm ord}\simeq A$. It has the least possible injection-cardinality
of any prime-ideal-free nontrivial Boolean algebra. In particular, for
every set $X$ in this model,
\[
 \begin{split}
 X\text{ is well-orderable}\quad\Longleftrightarrow\quad&
 \text{every nontrivial Boolean algebra }B\preccurlyeq X\\[-2pt]
 &\text{has a prime ideal}.
 \end{split}
\]

\textbf{Finite-cone realization and an unavoidable new subset.}
Put $S=[A]^{<\omega}$, $C_F=\{E\in S:F\subseteq E\}$, and
\[
 \mathcal I=\{D\subseteq S:(\exists F\in S)\ D\cap C_F=\varnothing\}.
\]
Canonically linearly extend each finite native order on $E$ by repeatedly
taking its lexicographically least remaining minimal point; call the
result $L_E$. The old sets
\[
 D_{ab}=\{E\in S:a,b\in E\text{ and }a<_{L_E}b\}
\]
satisfy the order constraints modulo $\mathcal I$. The subalgebra
generated by their classes inside $\mathcal P(S)/\mathcal I$ is another
size-$A$ algebra with no prime ideal. The entire quotient need not
have size $A$.

Even an ultrafilter on the \emph{old} algebra $\mathcal P(S)^W$ containing
all $C_F$ defines a native linear extension by
$a<_U b\Longleftrightarrow D_{ab}\in U$. Thus forcing with the positive
part of $\mathcal P(S)^W/\mathcal I$ necessarily adds a new subset of
$A^2$, and therefore of $A$. Moreover, the $A^2$-indexed dense open
sets deciding the Boolean values $[D_{ab}]$ have empty intersection.
In the Boolean completion they witness failure of binary distributivity
indexed by $A$. No claim about countable distributivity follows.
Genericity decides old subsets; it does not automatically give an
ultrafilter on the enlarged final powerset.

\textbf{Graph compactness.}
There is a graph $G$ with $|G|=|A|$ such that every finite subgraph is
$3$-colorable but $\chi(G)=4$. The order theory above is a finitely
satisfiable $3$-CNF theory after repeating literals. Encode Boolean
variables by opposite vertices adjacent to a fixed base vertex in a
triangle of colors $T,F,B$. A clause $x\vee y\vee z$ is equivalent to
\[
 \exists w\;\bigl(\operatorname{NAE}(x,y,w)
 \wedge\operatorname{NAE}(\neg w,z,F)\bigr),
\]
where NAE means ``not all equal.'' Each NAE constraint is represented
by a fresh triangle whose three vertices are respectively adjacent
to the three Boolean terminals. It has an extending $3$-coloring
exactly when NAE holds. A global $3$-coloring would solve the impossible
order theory. Every finite subgraph is colorable by finite
satisfiability, while a canonical $4$-coloring assigns Boolean values
uniformly and colors each gadget greedily. Tagged finite powers of
$A$ give the exact vertex cardinality.

Well-orderable graphs satisfy finite-color compactness in ZF by
transfinite recursion, choosing at each vertex the first color
preserving finite satisfiability. Thus $A$ is again a least possible
counterexample cardinal. In contrast, $2$-color compactness holds in
$W$: finite bipartiteness gives a canonical pair of parity classes in
each nontrivial connected component, and $\mathsf{AC}_2$ selects one
class per component. Here $\mathsf{AC}_2$ follows from OP.
For each fixed finite $k\geq3$, joining this graph with a complete graph on $k-3$ fresh vertices gives a graph whose finite subgraphs are $k$-colorable while its chromatic number is $k+1$. Thus finite-$k$ compactness fails for every such $k$.
The general compactness connection is classical; see de Bruijn--Erd\H{o}s
(1951) and Cowen--Hechler--Mihok (2002).
The explicit gadgets and infinite argument are checked in
\path{EXPLICIT_GRAPH_COMPACTNESS_INDEPENDENT_AUDIT.txt}; the Boolean
proof is in
\path{ORDER_OBSTRUCTION_BOOLEAN_ALGEBRA_AND_FINE_FORCING.txt}.

\subsection{BPI, small violations of choice, and limits of compactness}

Under $\mathsf{SVC}(A)$, BPI is equivalent to the existence of a fine
ultrafilter on $[A]^{<\omega}$, equivalently a prime ideal in
$\mathcal P([A]^{<\omega})/\mathcal I$. This is the Pincus--Blass local
reflection recorded in Ryan-Smith, \emph{Local reflections of choice},
Proposition 3.13. The preceding theorem settles this particular local
test negatively in $W$; it does not prove the abstract implication
$\mathsf{SVC}(A)+\mathsf{AC}_{\rm WO}+Q(A)+\mathsf{BPI}\Rightarrow
\mathsf{AC}$.

The model gives relative nonimplications
\[
 \mathsf{PP}\not\Rightarrow\mathsf{BPI},\qquad
 \mathsf{NDS}\not\Rightarrow\mathsf{BPI},\qquad
 \mathsf{FB}\not\Rightarrow\mathsf{BPI},
\]
and likewise with OEP as conclusion. In the reverse direction, the
classical basic Cohen model satisfies BPI and SVC and has an infinite
Dedekind-finite set. It fails PP, NDS, FB, and Ramsey's principle on
arbitrary infinite sets. This is an old model, not a new construction;
the SVC generator need not be its Cohen set. Tachtsis's 2024 Theorem
6.4 supplies the stronger opposite-direction example
$\mathsf{BPI}+\mathsf{AC}_\omega+\neg\mathsf{NDS}$, hence also
$\neg\mathsf{FB}$ and ordinary infinite homogeneity. Its exact
attribution was verified through Blass--Kulshreshtha's primary
manuscript, Theorem 10.5; the full Tachtsis article was not obtained.
These statements do not decide whether PP implies NDS in arbitrary
ZF models, or whether PP and BPI can coexist with $\neg\mathsf{AC}$.

Several useful subsidiary facts survived the unsuccessful positive-BPI
routes. FB gives every proper filter a member of minimum injection
cardinality; restricting an ultrafilter to that member makes it
cardinal-uniform. This is not a cofinal filter base or an ultrafilter
existence theorem. The native lower cones generate a proper filter
closed under all internal well-ordered intersections and avoiding every
well-orderable subset of $A$. An ultrafilter extending it need not retain
that completeness. In the original model every nonprincipal ultrafilter
on $A\subseteq\mathcal P(\omega_1)$ is countably incomplete: first
push forward by the first disagreement with its coordinatewise
ultrafilter decisions, obtaining a nonprincipal image on $\omega_1$;
then use the ground injection $\omega_1\to\mathcal P(\omega)$ and
repeat, obtaining a nonprincipal image on $\omega$.

Likewise, BPI and PP give a valid small-carrier Henkin lemma:
a finitely satisfiable finitary theory whose symbols inject into $A$
has a model of cardinality at most $A$, provided $A^{<\omega}\simeq A$.
This does not transport the original monoid action through an arbitrary
bijection of carriers. Finite-factor ultraproduct arguments also do
not force strict external cardinal descent: in ZFC every unbounded
finite product $\prod_n k(n)/U$, for nonprincipal $U$ on $\omega$,
has size $2^{\aleph_0}$, even after finitely many iterated logarithms
of the factors. These boundaries are recorded in
\path{STRUCTURAL_BPI_PHASE1.txt},
\path{ANCESTOR_FILTER_COMPACTNESS_AND_ULTRAFILTER_INCOMPLETENESS.txt},
\path{HENKIN_SMALL_CARRIER_AND_STRUCTURE_OBSTRUCTION.txt}, and
\path{ULTRAPRODUCT_AND_QUOTIENT_OBSTRUCTIONS.txt}.

\subsection{Parameter changes and preservation}

\textbf{Audited parameter theorem.}
Let the ground satisfy ZFC. Suppose $\kappa$ is regular uncountable,
$\mu$ is regular, and
\[
 2^{<\kappa}<\mu,\qquad
 \lambda^{<\kappa}<\mu\quad
 \text{for every infinite ground cardinal }\lambda<\mu.
\]
Use the source monoid of size $\mu$, $|I|=\mu^+$, Cohen conditions of
size $<\kappa$, and diagram supports of size at most $\mu$. The audited
adaptation yields the same SVC, ordinal-choice and quotient package,
hence PP, FB, NDS and $\neg\mathsf{AC}$. The quotient proof requires
an explicit repair: bounded arrays indexed below $\omega_1$ become
arrays of size $<\kappa$, controlled by the strengthened no-new-ground-
sequences lemma. The countable germ construction remains countable.
The native six-cycle and real-height obstruction remain valid, giving
$\neg\mathsf{BPI}$ and $\neg\mathsf{DC}_{\omega_1}$ in every variant,
although ordinary DC holds.

Two internal thresholds are distinct. If $W_\lambda$ says that every
set is injection-comparable with the well-ordered cardinal $\lambda$,
then
\[
 w(A)=\mu,\quad h(A)=\mu^+,\qquad
 W_\lambda\ \Longleftrightarrow\ \lambda\leq\mu,
 \qquad
 \mathcal P(\alpha)\text{ is well-orderable}\
 \Longleftrightarrow\alpha<\kappa.
\]
Here $w(A)$ denotes the largest well-ordered cardinal injecting into
$A$. For $\alpha<\kappa$, the entire powerset equals its ground
counterpart. Under $V=L$, write $M_{mn}$ for $\kappa=\aleph_m$,
$\mu=\aleph_n$, $1\leq m<n<\omega$. All ground alephs are preserved:
\begin{center}
\begin{tabular}{c|cc|cc}
Model & $\kappa$ & $\mu$ & $W_{\aleph_3}$ &
$\mathcal P(\aleph_1)$ well-orderable\\\hline
$M_{12}$ & $\aleph_1$ & $\aleph_2$ & no & no\\
$M_{13}$ & $\aleph_1$ & $\aleph_3$ & yes & no\\
$M_{23}$ & $\aleph_2$ & $\aleph_3$ & yes & yes
\end{tabular}
\end{center}
These are specific separations over a common strong package, not
unrestricted independent assignments of two sentences. The consistency
interpretation uses the source's formula-by-formula forcing method over
$L$; ordinary $\operatorname{Con}(\mathsf{ZF})$ suffices, with no
inaccessible cardinal. The complete adaptation is not Lean-formalized.
See \path{LARGER_MU_SECOND_INDEPENDENT_AUDIT.txt}.

Without GCH, put $\theta=(2^{<\kappa})^V$. The ground Cohen forcing is
$\theta^+$-cc and preserves $\kappa$ and its regularity. Every old
binary string of length $<\kappa$ occurs as a block of any one generic
row. The row's block map therefore collapses the old set of such strings,
of size $\theta$, to cardinality $\kappa$ internally in $W$. Thus
\[
 (\kappa^+)^W=(\theta^+)^V,\qquad
 (\kappa^{<\kappa})^W=\kappa.
\]
Ordinal-function locality transfers the relevant preservation from
the Cohen extension, not from the full outer product extension.
The exact no-GCH pair-coloring argument respects this identification;
see \path{NO_GCH_HOMOGENEOUS_BOUND_RECHECK.txt}.

\textbf{Preservation theorem.}
If $W\subseteq N$ have the same ordinals, $N$ adds no subsets of $A$,
and $N$ satisfies $\mathsf{SVC}(A)$, then $N$ retains
$\mathsf{AC}_{\rm WO}+Q(A)+\neg\mathsf{AC}$ and its PP/FB consequences.
Indeed, $A^2\simeq A$ makes every new map on $A$ have an old kernel.
An old non-well-orderable quotient cannot become well-orderable without
well-ordering $A$. Old ordinal-choice then supplies representatives for
kernels of new ordinal-valued maps. SVC reduces general ordinal-choice
to these maps. Ordinary set forcing always preserves SVC by evaluating
a ground presentation of the set of subnames.

This criterion has concrete instances. If the outer ZFC forcing has
size $\chi$, take regular $\lambda>\chi$ and the ground forcing
$\operatorname{Add}(\lambda,\delta)^V$. Easton's small-forcing argument
makes it $<\lambda$-distributive in the outer extension. The audited
intersection lemma $W[K]\cap V[F][H]=W$ then gives no new subsets of
$A^2$ when $\lambda$ exceeds its outer size. Taking $\lambda$ also above
the outer size of $\mathcal P([A]^{<\omega})^W$ preserves the BPI truth
value; here that means preserving its failure. For the original
$L$-ground model one may use $\lambda=\aleph_5^V$ and
$\delta\geq\lambda^+$ with $\delta^\lambda=\delta$. The largest
well-ordered cardinal injecting into $\mathcal P(\lambda)$ becomes
$\delta$, while that powerset stays non-well-orderable. These are
arithmetic variations, not further weak-choice separations by themselves.
See \path{HIGH_FORCING_PRESERVATION_INDEPENDENT_AUDIT.txt}.

\subsection{Exact pair and finite-arity Ramsey spectra}

For an infinite well-ordered cardinal $\lambda$, use the temporary
notation $\mathrm{NR}^{r,\mathrm{fin}}_\lambda$ for: every finite
coloring of $[X]^r$, for every non-well-orderable $X$, has a homogeneous
subset of cardinality $\lambda$. The requirement concerns all sets and
colorings, not only a distinguished witness. Under AC it is vacuous.
Ordinary Ramsey homogeneity with an unspecified infinite size already
follows from DC here.

\textbf{Pairs, without GCH.}
Throughout the audited family, the largest $\lambda$ for binary pair
colorings is exactly $\kappa$. More strongly, every pair coloring with
an ordinal color cardinal $1\leq\eta<\kappa$ on a non-well-orderable
set has a homogeneous subset of size $\kappa$. At $\eta=\kappa$,
the coloring
\[
 c(\{a,b\})=\min(a\mathbin{\triangle}b),\qquad a,b\in A,
\]
has no homogeneous triple: three binary bits cannot be pairwise
different at the same first-difference coordinate. Thus $\kappa$ is
the exact least well-ordered color cardinal permitting a failure of
infinite pair homogeneity on a non-well-orderable set.

For the positive direction, encode a pulled-back coloring on
$\kappa^+$, the enumeration of $\kappa^{<\kappa}$, and ordinal-sized
injection witnesses into a set $C$ of ordinals. The inner model $L[C]$
computes the required cardinals and $2^{<\kappa}=\kappa$ correctly.
Apply the Erd\H{o}s--Rado theorem there, in the form
$(2^{<\kappa})^+\to(\kappa+1)^2_\eta$ for $\eta<\kappa$.
The resulting homogeneous set belongs to $W$. Use
$\kappa^+\preccurlyeq A$ and then the copy of $A$ in every
non-well-orderable set. The precise primary source is
Baumgartner--Hajnal--Todorcevic, \emph{Extensions of the Erd\H{o}s--Rado
Theorem}, Theorem 2.1, p.~4.

For the negative direction, color $\{a,b\}$ by the parity of their
first difference. Fresh-copy locality first makes every homogeneous
subset well-orderable. A Cohen $\Delta$-system argument excludes size
$\kappa^+$: thin conditions to equal row patterns and then prescribe
a fresh first difference of the opposite parity. Both colors have
homogeneous subsets of size $\kappa$, supplied by rows whose first $1$
occurs at increasing coordinates of a fixed parity. In the no-GCH case
the argument uses the preserved ground $\theta^+$, which is the
internal $\kappa^+$. Consequently
\[
 \mathrm{NR}^{2,\mathrm{fin}}_\lambda
 \quad\Longleftrightarrow\quad \lambda\leq\kappa.
\]
Sources and complete proofs are
\path{EXACT_HOMOGENEOUS_CARDINAL_BOUND.txt} and
\path{GLOBAL_NON_WELL_ORDERABLE_RAMSEY_THRESHOLD.txt}.

\textbf{Higher arities over $L$.}
In $M_{mn}$ the exact threshold for all finite color counts is
\[
 \boxed{\quad
 \max\{\lambda:\mathrm{NR}^{r,\mathrm{fin}}_\lambda\}
       =\aleph_{\max(m-r+2,0)},\qquad 2\leq r<\omega.
 \quad}
\]
\begin{center}
\begin{tabular}{c|cccc}
Model & pairs & triples & $4$-tuples & $5$-tuples\\\hline
$M_{12},M_{13}$ & $\aleph_1$ & $\aleph_0$ & $\aleph_0$ & $\aleph_0$\\
$M_{23}$ & $\aleph_2$ & $\aleph_1$ & $\aleph_0$ & $\aleph_0$\\
$M_{34}$ & $\aleph_3$ & $\aleph_2$ & $\aleph_1$ & $\aleph_0$
\end{tabular}
\end{center}
The positive proof applies finite-arity Erd\H{o}s--Rado inside $L[C]$;
all powersets needed for its finite GCH calculation lie strictly below
$\kappa$ and are unchanged. The negative proof begins with Sierpi\'nski's
comparison of a ground ordinal enumeration and binary lexicographic
order, then steps up using consecutive first differences. A tuple is
colored by the preceding color together with an increasing/decreasing
tag, or by an additional nonmonotone color. Finite Ramsey excludes an
infinite homogeneous set of the latter color. DC makes a homogeneous
set in either remaining orientation well-ordered by lex or reverse lex;
its first differences yield a forbidden homogeneous set for the prior
coloring. This avoids an auxiliary well-order of $\mathcal P(\kappa)$.
The negative triple witness uses four colors; successive constructions
use at most $5\cdot2^{r-3}-1$ colors, not claimed optimal.
See \path{HIGHER_ARITY_ROOT_INDEPENDENT_AUDIT.txt} and Erd\H{o}s--Rado
(1956), Theorem 39, pp.~467--469.

\textbf{Binary triples remain a separate question.}
The saved bounds are
\[
\begin{array}{c|ccc}
 &\text{proved positive}&\text{unresolved here}&\text{proved negative}\\\hline
 M_{12},M_{13}&\omega&\omega_1&\omega_2\\
 M_{23}&\omega_1&\omega_2&\omega_3
\end{array}
\]
for homogeneous cardinalities of arbitrary binary triple colorings on
non-well-orderable domains. The binary orientation coloring of adjacent
first differences has exact largest homogeneous cardinality $\kappa$,
in both colors; this gives the stated negative upper bounds, not a
universal positive theorem at $\kappa$. Every one of the sixteen maps
merging the particular four-color triple witness into two colors has
an $\omega_1$-sized homogeneous set in $M_{12}$ and $M_{13}$. Projections
use row patterns; the XOR case uses a countable binary tree of
coordinates labeled by nested real intervals. Thus simple color merging
cannot settle the missing negative result. The classical binary
stepping-up proof uses an auxiliary well-order unavailable internally.
See \path{BINARY_TRIPLE_ROOT_INDEPENDENT_AUDIT.txt}. No claim is made
that the unresolved sentences are open in the entire literature.

\subsection{No-go results for repairing the model}

Increasing $|I|$ alone cannot produce BPI while retaining the absorption
probes, all valid three-copy mergers, root-fixing transport, local
decision, and Cohen conditions mentioning fewer than $\mu$ labels.
Those ingredients already prove the six-cycle obstruction. More
abstractly, the literal infinite-product version of finite Ramsey
homogeneity is false: parity relative to representatives modulo finite
coordinate disagreement gives a binary coloring with no monochromatic
rectangle having infinitely many nonsingleton factors.

Two concrete order-based changes were audited. Replacing the full
coordinate group by automorphisms of a saturated dense linear order
adds a supported cut map $A\twoheadrightarrow\mu^+$, while support
locality still forbids $\mu^+\hookrightarrow A$. Thus PP and
$\mathsf{AC}_{\rm WO}$ fail. Alternatively, decorate every entire
diagram with a ground linear extension of $T$, preserving that order
under restriction. The root-pair amalgamation survives, but the
internal profile map surjects onto all ground linear extensions of
$T$, a set of ground size $2^\mu$. Again PP and ordinal-choice fail.
This differs from finite order forcing over $W$, discussed elsewhere
in this report, whose orbit profiles are new and need not be ground.
See \path{ORDERED_SUBGROUP_CUT_QUOTIENT_OBSTRUCTION.txt} and
\path{ORDER_REPAIR_ROOT_AUDIT.txt}.

For restricted palettes $\mathcal C$ of orders on $T$, closed under
$\vartheta\mapsto\vartheta_t$ with
$x\mathrel{\vartheta_t}y\Longleftrightarrow tx\mathrel\vartheta ty$,
there is no nonempty finite palette, even without requiring extension
of native divisibility. Finite Ramsey produces four probes monotone
up to reversal in every palette order; a five-chart tournament merger
would give a finite set with no least element. For infinite
$|\mathcal C|=\lambda$, the same argument and Erd\H{o}s--Rado give the
necessary bound $\mu\leq2^{2^\lambda}$. It is not an optimal bound.
Palettes of size at most $\mu$ can satisfy pullback closure and every
acyclic finite request. A universal lift can also be constructed for
one prescribed palette, using cross-denominator probe cones sharing
exactly the root. The unresolved issue is closure of one final palette
under universal lifts for all its own witness types, followed by support
descent and quotient classification. No positive-BPI model or universal
impossibility theorem results. See \path{ORDER_PALETTE_ROOT_AUDIT.txt}.

Ransom's 2026 dissertation gives an exact forcing-aware criterion for
BPI (Definition 4.3.1 and Theorem 4.3.3, pp.~93--99), involving dense
sets, bases for the normal filter, orbit colorings, and compatibility
below all common extensions of finitely many translated conditions.
The criterion concerns the full Cohen--diagram product. A Ramsey
property of the abstract monoid alone does not verify it. Its
finite-support examples do not automatically cover these $\mu$-sized
supports. The six-cycle theorem now rules out a successful verification
for the unchanged model. The earlier precise compatibility obligations
remain documented in \path{SYMMETRIC_FILTER_PHASE3.txt} and
\path{SYMMETRIC_FILTER_PHASE4_LARGE_I.txt}.

\subsection{Source corrections, attribution, and remaining boundaries}

The implication $\mathsf{AC}_{\rm WO}\Rightarrow\mathsf{DC}$ is valid
in pure ZF: Jech, Theorem 8.2, pp.~121--122, attributes it to Jensen.
An early apparent contrary table entry was an erroneous HTML
reconstruction, corrected by inspecting the actual primary PDF and
its prose. Here ordinal-choice means choice for well-orderable
\emph{index families}, not arbitrary families of well-orderable
members. The verified chain is
$\mathsf{PP}\Rightarrow\mathsf{AC}_{\rm WO}\Rightarrow\mathsf{DC}$;
it does not yield $\mathsf{DC}_{\omega_1}$.

A different source conflict remains recorded rather than resolved by
recency. Pincus's original announcement and publisher summary, Ransom's
2025 paper, and Asper\'o--Karagila's 2020 manuscript explicitly attribute
models of $\mathsf{ZF}+\mathsf{DC}+\mathsf{BPI}+\neg\mathsf{AC}$ to
Pincus. Kowalczyk--Zapletal, arXiv:2607.16747v1, instead calls the
corresponding implication to AC open; its actual theorem establishes
a restricted compactness fragment. The complete relevant Pincus proof
was not retrieved in this investigation. This discrepancy is not used
as a premise, and no claim that a positive BPI variant would solve
that historical problem is made. See
\path{DC_BPI_SOURCE_CONFLICT.txt}.

The bounded novelty check found no BPI or order-extension statement in
the twelve text files of the pinned PP manuscript and no prior assertion
of this particular PP/FB/NDS countermodel to BPI in the sources checked.
That is limited search evidence, not a priority certificate. The
classical compactness implications, symmetry-cycle method, Ramsey
theorems and coloring constructions retain their established
attributions. The source construction, not this report, proves the
foundational PP consistency theorem. The additional results here are
model consequences and explicitly audited adaptations.

\paragraph{Primary references used in this section.}
Stable source links are:
\href{https://gwern.net/doc/math/1973-jech-theaxiomofchoice.pdf}{Jech (1973)};
\href{https://doi.org/10.2307/2586759}{Felgner--Truss (1999), OEP and BPI};
\href{https://academic.calliope.mx/files/papers/2412.13785.pdf}{Ryan-Smith, Proposition 3.13};
\href{https://arxiv.org/abs/2511.21684v1}{Ransom (2025)} and
\href{https://escholarship.org/uc/item/0b35t7sd}{Ransom's 2026 dissertation};
\href{https://doi.org/10.1215/00294527-2024-0020}{Tachtsis (2024)} and
\href{https://arxiv.org/abs/2310.09643v3}{Blass--Kulshreshtha, version 3};
\href{https://arxiv.org/pdf/math/9311207}{Baumgartner--Hajnal--Todorcevic,
Theorem 2.1};
\href{https://www.renyi.hu/~p_erdos/1956-02.pdf}{Erd\H{o}s--Rado (1956)};
\href{https://users.renyi.hu/~p_erdos/1965-14.pdf}{Erd\H{o}s--Hajnal--Rado
(1965), Lemma 5A}; and
\href{https://www.jams.jp/notice/scmj/56-2.html}{Cowen--Hechler--Mihok (2002)}.
All named local audit files belong to the accompanying
\path{evidence/original_model/} archive directory (mapped from the original \path{BPI_RPP_five_hour_research/} record). Their original
intermediate statuses are preserved; the later independent audits and
this consolidation specify which claims supersede them.
